% Experiments designed/concieved by Vijay Erramilli. Code written by Vijay Erramilli and Codex

\subsection{Qwen2.5-7B: the pattern transfers to a larger model}
\label{sec:qwen-induction}

We test that transfer on the base Qwen2.5-7B model
\cite{qwen2024qwen25} using an induction-copy task with an exact behavioral
target and a different set of heads to intervene on
\cite{olsson2022induction}. The result applies only to prompts that the
unedited model already solves with a clear margin.

Each prompt first contains three key--value pairs and later repeats one key.
For example, after ``A red, B blue, C green, ... B'', the desired continuation
among the three values is ``blue.'' Red and green are the declared distractors.
For each prompt, we score how strongly the model favors the target over the two
distractors combined. For prompt $i$, define this margin as
\[
M_i=\operatorname{logit}_i(\text{target})-
\operatorname{logmeanexp}
\bigl(\operatorname{logit}_i(\text{distractor}_1),
      \operatorname{logit}_i(\text{distractor}_2)\bigr),
\]
where
$\operatorname{logmeanexp}(a,b)=\operatorname{logsumexp}(a,b)-\log 2$.
The effect of an ablation mask is the drop in that margin:
\[
Y_i(m)=M_i(\varnothing)-M_i(m).
\]
A positive $Y_i(m)$ means that ablating those heads weakens the target relative
to the two alternatives. Four prompt families combine sequence lengths 32 and
64 with repeat gaps 8 and 16.

\subsubsection{Define eligible prompts before discovering heads}

Copy-v1 required 95\% baseline accuracy among the three candidates. Qwen reached
92.97\%, so the run stopped before any attention or intervention result was
opened. Copy-v2 did not lower the threshold. Instead, before analysis it defined
an eligible prompt as one for which Qwen chooses the target among the three
candidates, has margin $M_i\geq\log 4$, and produces finite scores. Discovery
prompts must also agree under the eager and SDPA attention implementations.
Both protocols recorded their source and configuration in a private commit
before model outcomes. This records data-access order but is not an independent
preregistration.

Before scoring, we assign 3,072 fresh prompts to separate data banks. A total
of 2,612 prompts (85.03\%) pass the eligibility and coverage checks. Hash order,
not model score, then selects 1,536 prompts for reference, discovery, head
fitting, confirmation, calibration, and final held-out testing. These roles
remain separate throughout the experiment.

On the unfiltered prompts, Qwen produces the desired next token only 24.45\% of
the time. The task therefore asks a bounded question: when Qwen already
distinguishes the planted value from the two distractors, which heads support
that distinction, and how well can we predict their intervention effects? It
does not test unconstrained generation.

\subsubsection{Held-out interventions confirm the eight-head panel}

Attention can suggest heads, but only an intervention shows that they affect
the output. We therefore use four frozen steps:

\begin{enumerate}
  \item Rank all 784 query heads by attention to the target value relative to
  the distractors.
  \item Apply individual mean ablations to only the top 32 on a separate
  head-fitting bank.
  \item Select at most two heads per layer using a lower-bootstrap-bound rule.
  This yields L22H3, L22H4, L26H15, L14H0, L14H6, L16H13, L17H18, and L23H13,
  spanning six layers.
  \item Test the selected panel against matched control heads on a separate
  confirmation bank.
\end{enumerate}

Each selected head receives a low-induction control from the same layer and
grouped-query-attention key--value group. Controls are matched using attention
specificity and output-norm distance, not causal outcomes. On 256 separate
confirmation prompts, the selected heads lower the candidate margin much more
than matched controls, with the paired interval excluding zero; a frozen
zero-ablation check gives the same ordering. Table~\ref{tab:qwen-copy-v2}
reports the estimates. Thus the held-out intervention results, not the attention
scores by themselves, justify including the panel in the benchmark.

\subsubsection{Pair terms and direct loss improve held-out performance}

Only after confirmation passes do we freeze all $2^8=256$ combinations of the
eight heads. One mask might remove heads 1 and 3; another might remove 2, 4,
and 8. We split this complete mask set into 128 calibration masks and 128
held-out masks. Calibration effects average 256 prompts; held-out effects
average a separate 512. The held-out prompts form 128 four-family clusters, so
every mask is tested at every sequence-length and repeat-gap condition.

The additive observer assumes that removing heads 1 and 3 has the sum of their
separate effects. The pairwise observer, a quadratic model for these binary
masks, adds all 28 head pairs and can represent a head whose effect changes
after another head is removed. At the primary budget of 128 measurements, the
pairwise model improves held-out-mask MAE over the additive observer, with the
paired interval excluding zero (Table~\ref{tab:qwen-copy-v2}). The pairwise
model also improves MAE at every smaller nested budget; the 16-measurement fit
has more coefficients than measurements and is therefore descriptive.

Every action selector uses the same 37-column pairwise basis, the same three
targets, and the same 16 frozen action pools. Each pool contains eight nonempty
masks and the exact no-op. The observers differ only in
what they predict:

\begin{itemize}
  \item \textbf{Natural mean} predicts $\mathbb E_i[Y_i(m)]$.
  \item \textbf{Transformed mean} predicts the distance between that mean and
  the target.
  \item \textbf{Direct loss} predicts the action loss itself,
  $\mathbb E_i|Y_i(m)-t|$.
\end{itemize}

The targets are 0.25, 0.50, and 0.75 of the confirmed full-panel effect. We
freeze every prediction and action before opening held-out effects.

\begin{table}[H]
\centering
\scriptsize
\setlength{\tabcolsep}{3.5pt}
\begin{tabular}{@{}llrrl@{}}
\toprule
Test & Comparison & Reference & Result & Improvement [95\% CI] \\
\midrule
Causal & Mean replacement & Control $0.170$ & Selected $3.454$
  & $3.284\ [3.053,3.520]$ \\
Robustness & Zero ablation & Control $0.528$ & Selected $4.080$
  & $3.552\ [3.306,3.801]$ \\
\midrule
Prediction & 16 measurements$^\dagger$ & Additive $0.354$ & Pairwise $0.100$
  & $0.254\ [0.182,0.295]$ \\
& 40 measurements & Additive $0.202$ & Pairwise $0.045$
  & $0.157\ [0.103,0.173]$ \\
& 64 measurements & Additive $0.157$ & Pairwise $0.039$
  & $0.117\ [0.067,0.134]$ \\
& 128 measurements & Additive $0.121$ & Pairwise $0.040$
  & $0.081\ [0.025,0.101]$ \\
\midrule
Action loss & Natural mean & $0.955$ & Direct loss $0.870$
  & $0.085\ [0.053,0.118]$ \\
& Transformed mean & $0.995$ & Direct loss $0.870$
  & $0.125\ [0.083,0.173]$ \\
& Exact no-op & $1.727$ & Direct loss $0.870$
  & $0.857\ [0.805,0.908]$ \\
\bottomrule
\end{tabular}
\caption{Qwen2.5-7B Copy-v2 results. Causal entries report mean
candidate-margin drop on 256 confirmation prompts. Prediction entries report MAE on 128
held-out masks; positive improvement is additive minus pairwise MAE. Action
entries average absolute target loss over three targets and 16 pools; positive
improvement is comparator loss minus direct loss. Intervals use 5,000 paired
percentile bootstrap draws over the appropriate prompt clusters, masks, and
pools. The dagger marks the underdetermined, descriptive 16-measurement fit.
The 128-measurement prediction row and aggregate action rows are the
prespecified primary tests.}
\label{tab:qwen-copy-v2}
\end{table}

Direct loss lowers aggregate loss by 8.9\% relative to natural mean, 12.6\%
relative to transformed mean, and 49.6\% relative to no-op. These numbers give
equal weight to the three targets. At the largest target, direct loss and
natural mean tie within uncertainty: natural minus direct loss is $-0.008$,
with interval $[-0.038,0.021]$. Appendix
Table~\ref{tab:qwen-per-target} reports each target separately.

On 512 matched non-induction continuations, a secondary check finds no larger
clean-to-intervened KL for direct loss than for the natural-mean and
transformed-mean controls ($0.0115$ versus $0.0157$ and $0.0160$); this is not a
general safety result.

Across GPT-2 and this bounded Qwen task, pair terms improve held-out effect
prediction and observers trained on realized loss choose lower-loss
interventions in their frozen aggregates; the Qwen result, however, remains
limited to one pinned model, the eligible prompts, the confirmed eight-head
panel, and the declared masks and neither identifies a complete induction
circuit nor establishes general in-context learning.

These studies choose model interventions from predicted effects. What changes
when the action is instead a scarce safety response and different mistakes
carry different consequences? The next operational mode keeps the same
task-relative contract but evaluates observers under fixed intervention
budgets.
